\section*{Apendice E: Definiciones Matematicas de Metricas}

\subsection*{Frecuencia Cardiaca (HR)}

$$HR = \frac{\text{Numero de latidos}}{T_{\text{minutos}}} \quad [\text{bpm}]$$

Donde $T$ es la duracion del periodo de observacion en minutos.

\subsection*{Variabilidad de Frecuencia Cardiaca (HRV - SDNN)}

$$\text{HRV (SDNN)} = \sqrt{\frac{1}{N-1} \sum_{i=1}^{N} (NN_i - \overline{NN})^2} \quad [\text{ms}]$$

Donde:
\begin{itemize}
    \item $NN_i$ = duracion del i-esimo intervalo Normal-to-Normal (entre latidos R-R)
    \item $\overline{NN}$ = promedio de todos los intervalos NN
    \item $N$ = numero total de intervalos
\end{itemize}

HRV baja indica dominancia simpatica (estrés, fatiga).
HRV alta indica regulacion parasimpatica optima (recuperacion).

\subsection*{Tasa Respiratoria (BR)}

$$BR = \frac{\text{Numero de ciclos respiratorios}}{T_{\text{minutos}}} \quad [\text{brpm}]$$

\subsection*{Potencia EEG en Bels}

$$P_{\text{Bels}} = 10 \log_{10}\left(\frac{P_{\text{señal}}}{P_{\text{referencia}}}\right) \quad [\text{Bels}]$$

Escala logaritmica que comprime rangos dinamicos amplios.
Rango tipico en FatigueSet: -1 a +1 Bels (0.1 a 10 veces potencia de referencia).

\subsection*{Correlacion de Pearson}

$$r = \frac{\sum_{i=1}^{n} (X_i - \overline{X})(Y_i - \overline{Y})}{\sqrt{\sum_{i=1}^{n} (X_i - \overline{X})^2 \sum_{i=1}^{n} (Y_i - \overline{Y})^2}}$$

Rango: -1 (correlacion negativa perfecta) a +1 (correlacion positiva perfecta).
$r = 0$ indica ningun correlacion lineal.

\subsection*{Desviacion Estandar (SD)}

$$SD = \sqrt{\frac{1}{N-1} \sum_{i=1}^{N} (x_i - \overline{x})^2}$$

Medida de dispersion - cuanto varian los valores alrededor de la media.
